% =============================================================================
% RESUMEN / ABSTRACT
% =============================================================================
\newpage
\setcounter{page}{2}
\thispagestyle{fancy}

{\fontsize{16pt}{19pt}\selectfont\color{uaxblue} RESUMEN}\par

El acceso a historiales médicos electrónicos está severamente limitado por las normativas de privacidad vigentes, como el Reglamento General de Protección de Datos (GDPR). Esta imposibilidad de compartir información frena la innovación en el sector sanitario. Para solventar este bloqueo, este Trabajo de Fin de Máster propone un marco basado en Inteligencia Artificial Generativa diseñado para sintetizar gemelos de datos tabulares clínicos, garantizando el anonimato legal sin sacrificar el valor predictivo original.

A nivel metodológico, se ha diseñado un flujo de trabajo empleando un conjunto de datos médicos reales, caracterizado por un severo desbalanceo de clases y variables mixtas. Se han entrenado y evaluado diversas arquitecturas de estado del arte, haciendo especial hincapié en la adaptación de un modelo probabilístico de difusión al dominio tabular. El rendimiento del ecosistema ha sido auditado mediante métricas de fidelidad estadística y protocolos de validación cruzada para certificar la utilidad predictiva. Además, se han integrado auditorías empíricas rigurosas para cuantificar el riesgo de memorización de pacientes y la resistencia frente a ataques de inferencia de membresía.

Los resultados empíricos demuestran la viabilidad técnica de la arquitectura de difusión progresiva frente a otros paradigmas generativos. Las distribuciones generadas replican fielmente las fronteras de decisión de los datos originales, ofreciendo una paridad predictiva excepcional al entrenar algoritmos externos. Paralelamente, las pruebas de privacidad evidencian que se anulan las vulnerabilidades de reidentificación, situando la probabilidad de éxito de cualquier atacante en niveles de mero azar.

En conclusión, la metodología desarrollada certifica que es posible democratizar los datos sensibles para la investigación. Al satisfacer los estrictos requerimientos de anonimización exigidos por la normativa europea, el sistema permite transformar repositorios restringidos en recursos explotables, reduciendo drásticamente los bloqueos burocráticos y los tiempos de acceso a los datos en proyectos corporativos de inteligencia artificial.

{\fontsize{13pt}{16pt}\selectfont\color{uaxblue} Palabras clave}\par

Inteligencia Artificial Generativa, Datos Tabulares Sintéticos, Modelos de Difusión, Privacidad de Datos, GDPR.

{\fontsize{16pt}{19pt}\selectfont\color{uaxblue} ABSTRACT}\par

Access to electronic medical records is severely limited by current privacy regulations, such as the General Data Protection Regulation (GDPR). This inability to share information hampers innovation in the healthcare sector. To overcome this bottleneck, this Master's Thesis proposes a framework based on Generative Artificial Intelligence designed to synthesize digital twins of clinical tabular data, guaranteeing legal anonymity without sacrificing the original predictive value.

Methodologically, a workflow has been designed using a real medical dataset, characterized by severe class imbalance and mixed variables. Various state-of-the-art architectures have been trained and evaluated, with special emphasis on adapting a denoising diffusion probabilistic model to the tabular domain. The performance of the ecosystem has been audited using statistical fidelity metrics and cross-validation protocols to certify predictive utility. Furthermore, rigorous empirical audits have been integrated to quantify patient memorization risk and resilience against membership inference attacks.

Empirical results demonstrate the technical viability of the progressive diffusion architecture compared to other generative paradigms. The generated distributions faithfully replicate the decision boundaries of the original data, offering exceptional predictive parity when training external algorithms. Concurrently, privacy tests show that re-identification vulnerabilities are nullified, placing any attacker's probability of success at mere chance levels.

In conclusion, the developed methodology certifies that it is possible to democratize sensitive data for research. By satisfying the strict anonymization requirements demanded by European regulations, the system allows transforming restricted repositories into exploitable resources, drastically reducing bureaucratic roadblocks and data access times in artificial intelligence projects.

{\fontsize{13pt}{16pt}\selectfont\color{uaxblue} Keywords}\par

Generative Artificial Intelligence, Synthetic Tabular Data, Diffusion Models, Data Privacy, GDPR.

\vspace{0.5cm}
